% ==============================
% 使用 ctex 统一处理中文
% （自动加载 xeCJK，避免重复配置）
% ==============================
\usepackage[heading=false]{ctex}   % 关闭 ctex 自动重定义章节标题样式，便于出版社后期调整


% ==============================
% 字体配置（macOS 原生字体）
% XeLaTeX 专用
% ==============================
\usepackage{fontspec}              % XeLaTeX / LuaLaTeX 字体支持

\setmainfont{Times New Roman}      % 英文正文：出版最常用
\setsansfont{PingFang SC}          % 英文/中文无衬线：图注、标题
\setmonofont{Menlo}                % 等宽字体：代码、命令行

\setCJKmainfont{Songti SC}         % 中文正文：宋体，符合高教社习惯
\setCJKsansfont{PingFang SC}       % 中文黑体：标题、图表
\setCJKmonofont{PingFang SC}       % 中文等宽字体：苹方（macOS 自带）


% ==============================
% 处理 Unicode 上下标字符
% ==============================
\usepackage{newunicodechar}
\usepackage{siunitx}

% 只要在文档里看到这几个字符，就自动转为 LaTeX 标准上下标
\newunicodechar{⁰}{$^{0}$}
\newunicodechar{¹}{$^{1}$}
\newunicodechar{²}{$^{2}$}
\newunicodechar{³}{$^{3}$}
\newunicodechar{⁴}{$^{4}$}
\newunicodechar{⁵}{$^{5}$}
\newunicodechar{⁶}{$^{6}$}
\newunicodechar{⁷}{$^{7}$}
\newunicodechar{⁸}{$^{8}$}
\newunicodechar{⁹}{$^{9}$}
\newunicodechar{₀}{$_{0}$}
\newunicodechar{₁}{$_{1}$}
\newunicodechar{₂}{$_{2}$}
\newunicodechar{₃}{$_{3}$}
\newunicodechar{₄}{$_{4}$}
\newunicodechar{₅}{$_{5}$}
\newunicodechar{₆}{$_{6}$}
\newunicodechar{₇}{$_{7}$}
\newunicodechar{₈}{$_{8}$}
\newunicodechar{₉}{$_{9}$}
\newunicodechar{₊}{$_{+}$}
\newunicodechar{⁺}{$^{+}$}
\newunicodechar{⁻}{$^{-}$}
\newunicodechar{⁼}{$^{=}$}
\newunicodechar{⁽}{$^{(}$}
\newunicodechar{⁾}{$^{)}$}
\newunicodechar{・}{$\cdot$}
\newunicodechar{₋}{$_{-}$}
\newunicodechar{₌}{$_{=}$}
\newunicodechar{₍}{$_{(}$}
\newunicodechar{₎}{$_{)}$}
\newunicodechar{℃}{\si{\celsius}}
\newunicodechar{∶}{:}

% ==============================
% 页面版式（技术书籍紧凑风格）
% ==============================
\usepackage[
  a4paper,                         % A4 纸张
  left=2.5cm,                      % 左边距（装订侧）
  right=2cm,                       % 右边距
  top=2cm,                         % 上边距
  bottom=2cm                       % 下边距
]{geometry}


% ==============================
% 行距与段落
% ==============================
\usepackage{setspace}              % 行距控制
\setstretch{1.15}                  % 1.15 倍行距（技术书籍紧凑风格）

\usepackage{indentfirst}           % 章节后首段也缩进
\setlength{\parindent}{2em}        % 首行缩进 2 个汉字

% 段落间距
\setlength{\parskip}{0.2em}        % 段落之间的垂直间距，紧凑但仍可区分段落



% ==============================
% 数学公式与定理环境
% ==============================
\usepackage{amsmath}               % AMS 数学核心
\usepackage{amssymb}               % 数学符号
\usepackage{amsthm}                % 定理环境

\numberwithin{equation}{chapter}   % 公式编号按“章-序号”


% ==============================
% 图形、表格与浮动体
% ==============================
\usepackage{graphicx}              % 插入图片
\usepackage{caption}               % 图表标题控制
\usepackage{subcaption}            % 子图支持

\usepackage{booktabs}              % 三线表（出版社强烈推荐）


% ==============================
% 中文数学与中英文间距优化
% ==============================
\xeCJKsetup{
  CJKmath = true                   % 数学模式下允许中文
}


% ==============================
% 参考文献（GB/T 7714-2015）
% ==============================
\usepackage[
  backend=biber,                   % 使用 biber 引擎
  style=gb7714-2015                % 国标参考文献格式
]{biblatex}

\addbibresource{references.bib}    % 参考文献数据库文件

% ==============================
% PDF 页面插入支持
% ==============================
\usepackage{pdfpages}

